\section{Finite Weil arithmetic and the actual normalized theta periods}
\label{sec:finite-weil-arithmetic}

This section supplies the qualitative finite-Weil assertion retained in
TT13 and TT16.  It treats the actual representation and its actual
normalized vacuum.  In particular, the argument determines an algebraic
phase; preservation of a norm alone would determine only its absolute
square.  The consequences are the actual normalized jet field, and the
actual normalized definite-period field for every supplied bottom
triangle.  The existence of additional full adelic Weil representations
for new bottom spaces is separated at WA17--WA18.

Here are the common hypotheses for WA1--WA11.  Fix the source's CM host
$D$, standard additive character $\psi$ of sign $-1$, its specified
infinity embeddings $\theta$, Hermitian spaces $S,T$, self-dual full
$\mathcal O_E$-lattices $\Lambda_S,\Lambda_T$, and their specified Cartan
involutions at infinity.  Let $v$ be the actual unit vacuum
\texttt{IsVacO} for these data, and let $\omega$ be the actual
\texttt{IsWeilDatumO}, with its given splitting Hecke characters
$\chi_S,\chi_T$ satisfying \texttt{ResEta} at the respective ranks.
Thus Heisenberg covariance, continuity, norm preservation, rational
automorphy and the positive-scalar split normalization are hypotheses
on this particular representation.  We do not infer its existence
from finite arithmetic.  Write $\mathcal W=\operatorname{Hom}_E(S,T)$,
write $b$ for the source's oriented symplectic pairing, and put
\[
 U(x)f(w)=\psi\bigl(b(w,x)/2\bigr)f(w+x),\qquad
 \Lambda=\widehat{\mathfrak a}_D\,
          \operatorname{Hom}_{\widehat{\mathcal O}_E}
                 (\widehat\Lambda_S,\widehat\Lambda_T),
             \qquad\mathfrak a_D=D.\mathrm{host}.\mathrm{AA}.
                                                               \tag{WA.1}
\]
This is exactly the source's \texttt{latO}, including its fractional
ideal factor; the Hom lattice alone would generally be different.
Elements of $\Lambda$ have zero infinite component.  The identity
$\xi_0\mathfrak a_D\overline{\mathfrak a_D}
=\mathfrak D_{E/\mathbb Q}^{-1}$ and the perfect Hermitian
pairings on the two self-dual lattices show, by local dual bases,
that the finite annihilator of $\Lambda$ under $\psi\circ b$ is
exactly $\Lambda$.  This is the trace-dual calculation of VI2,
with the actual \texttt{latO} multiplier and present frames.
Throughout, $\overline{\mathbb Q}$ denotes the actual algebraic closure
of $\mathbb Q$ inside $\mathbb C$, and
\[
 \mathcal A_v=\operatorname{span}_{\overline{\mathbb Q}}
                     \{U(x)v:x_\infty=0\}.                  \tag{WA.2}
\]
Every span expression in this section is a finite sum.

\paragraph{WA1: finite Heisenberg phases and the algebraic span.}
If $x_\infty=y_\infty=0$, choose positive integers $a,b$ with
$ax,by\in\Lambda$.  Such integers exist because a finite adele is
integral away from finitely many places and each remaining lattice
denominator can be cleared by an integer.  Self-duality gives
$\psi(b(ax,by))=1$.  Hence $\psi(b(x,y))$ has finite order dividing
$ab$, and $\psi(b(x,y)/2)$ has finite order dividing $2ab$.
These numbers are algebraic, nonzero and of absolute value one.
The actual operator identity is
\[
 U(x)U(y)=\psi\bigl(b(x,y)/2\bigr)U(x+y),\qquad
 U(-x)U(x)=1.                                                 \tag{WA.3}
\]
Consequently $\mathcal A_v$ contains $v$, is closed under algebraic
linear combinations, and is preserved by every $U(x)$ with finite
$x$.  This uses the half-pairing in (WA.1); dropping its factor of
$1/2$ would change the phases at primes over $2$.

The literal lattice-vacuum law gives one uniquely determined sign
$\varepsilon(\lambda)\in\{1,-1\}$ with
$U(\lambda)v=\varepsilon(\lambda)v$.  Uniqueness uses $v\ne0$,
which follows from its norm one.  The sign is one on $2\Lambda$
and on either isotropic half of $\Lambda$ at a split place.
To obtain the sign for all of $\Lambda$ from these literal clauses,
use that multiplication by $2$ is invertible at every odd place,
while the host makes every place over $2$ split.  The two local
idempotent halves there, together with $2\Lambda$, generate
$\Lambda$.  Their individual operators fix $v$; the product
phases have square one by the integral self-pairing.  This proves
the claimed sign and shows that it factors through $\Lambda/2\Lambda$.
Equation (WA.3) gives the exact semicharacter relation for
$\varepsilon$.  No simultaneous sign choice for unrelated vacua
is used.

\paragraph{WA2: algebraic pairings and rigidity of a scalar.}
Use the source's covolume-one theta inner product, conjugate-linear
in its first variable.  It is positive on nonzero theta functions
and satisfies $\langle v,v\rangle=1$.  We have
\[
 \langle v,U(x)v\rangle=
 \begin{cases}\varepsilon(x),&x\in\Lambda,\\0,&x\notin\Lambda,
 \end{cases}\qquad x_\infty=0.                              \tag{WA.4}
\]
The first case is the vacuum law.  For the second, self-duality
provides $\lambda\in\Lambda$ with $\psi(b(\lambda,x))\ne1$.
Apply $U(\lambda)$ to both entries of the inner product and use
(WA.3).  The two signs cancel, leaving the original inner product
multiplied by $\psi(b(\lambda,x))$; it must therefore vanish.
Heisenberg invariance now expresses
$\langle U(x)v,U(y)v\rangle$ as a phase from WA1 times a value
of (WA.4).  Expansion of two finite sums proves
\[
               F,F'\in\mathcal A_v
       \quad\Longrightarrow\quad
               \langle F,F'\rangle\in\overline{\mathbb Q}.
                                                               \tag{WA.5}
\]
If $F\ne0$ and $F'=cF$, positivity and (WA.5) give
$c=\langle F,F'\rangle/\langle F,F\rangle\in\overline{\mathbb Q}$.
In particular an equality of actual vectors, rather than equality
of their norms, detects the scalar.

\paragraph{WA3: a finite Weil move has an algebraic projective expansion.}
For a finite pair $g=(g_S,g_T)$, there are $F\in\mathcal A_v$ and
$\kappa\ne0$ such that
\[
                         \omega(g)v=\kappa F,
                              \qquad F\ne0.                \tag{WA.6}
\]
At this stage no algebraicity of $\kappa$ is asserted.  We give
the finite-dimensional argument behind the source's signed-average
proof.

First recall the analytic uniqueness used here.  For a self-dual
finite lattice, impose its specified semicharacter law and the
Fock annihilation equations for a positive Cartan complex structure.
Choose a symplectic integral basis of the corresponding rational
trace lattice.  Restriction to a fundamental finite coset then
gives the usual entire theta function with that positive complex
structure.  Fourier expansion along one integral half determines
all coefficients from one coefficient: translating by the opposite
half gives the theta recurrence.  Positive definiteness makes the
resulting Gaussian series normally convergent.  Thus this space has
dimension at most one, just as in VI6.  The change of real coordinates
is used only for this analytic uniqueness; its coefficients need
not be algebraic.  A known nonzero vacuum makes the dimension one.

For completeness, this uniqueness also gives a finite algebraic
model at every deeper level.  Let $V_M$ be the Fock theta functions
fixed by $U(2M\lambda)$ for all $\lambda\in\Lambda$, where $M>0$.
On $V_M$ the operators
$T(\lambda)=\varepsilon(\lambda)U(\lambda)$ form a representation
of the finite group $\Lambda/2M\Lambda$.  The semicharacter identity
proves the representation law, and the factor $2$ makes its kernel
contain $2M\Lambda$.  The finite character projectors are
\[
  P_\alpha=\frac1{|\Lambda/2M\Lambda|}
       \sum_{\lambda\bmod 2M\Lambda}
                          \alpha(\lambda)^{-1}T(\lambda).   \tag{WA.7}
\]
They sum to the identity.  Perfect finite-lattice duality represents
each $\alpha$ by pairing with some
$w\in(2M)^{-1}\Lambda/\Lambda$.  Translating an $\alpha$-eigenvector
by $U(-w)$ gives the original lattice-vacuum law.  Analytic uniqueness
therefore says that each eigenspace is the line generated by the
corresponding $U(w)v$.  This proves that $V_M$ has a finite basis of
such translates and that all Heisenberg operators preserving this
level have algebraic matrices in that basis: they permute these
lines with the roots of unity of WA1 and the lattice signs.

The finite pair $g$ preserves the Fock equations.  One may justify
differentiation here without assuming a linear map on all functions
is continuous: norm preservation makes $\omega(g)$ an isometry on
the theta $L^2$ space; on a fixed smooth finite level, archimedean
Heisenberg difference quotients converge in $L^2$.  Covariance moves
these quotients through that isometry.  The smooth output has the
resulting derivatives, giving the same annihilation equations.
Choose $M$ so that $\omega(g)v\in V_M$ and
$g\Lambda\subset(2M)^{-1}\Lambda$; increase $M$ if needed to have
$M\Lambda\subset g\Lambda$.  Commensurability gives all these
conditions.  The equations
\[
             U(g\lambda)f=\varepsilon(\lambda)f
                         \quad(\lambda\in\Lambda)          \tag{WA.8}
\]
are algebraic linear equations in the finite basis of $V_M$.
Their solution space is the one-dimensional vacuum line for the
moved self-dual lattice, by the same uniqueness argument.  An
intersection of kernels in a finite-dimensional space is determined
by finitely many of its equations.  Row reduction over
$\overline{\mathbb Q}$ therefore supplies a nonzero algebraic
coordinate vector $F$ spanning it.  This is (WA.6).

Equivalently, the source sums $\varepsilon(u)U(u)\phi$ over
$\Lambda/2M\Lambda$.  Its inner product with $v$ is the number of
summands times $\langle v,\phi\rangle$.  A nonzero character
projection supplies a translate with nonzero overlap, and applying
$\omega(g)$ to that average gives (WA.6).  Both proofs retain the
same unknown scalar.

\paragraph{WA4: a principal finite congruence fixes the actual vacuum.}
The simultaneous lattice stabilizer preserves the vacuum line:
covariance preserves the lattice law and its signs after a
sufficiently deep congruence, and a finite move preserves the
Fock equations as in WA3.  Shrink first to a compact open subgroup
on which the signs are preserved.  Thus
$\omega(k)v=c(k)v$ there.  The map $c$ is a character, has absolute
value one, and is continuous: evaluate at any point where $v$ is
nonzero and use the literal continuity clause of
\texttt{IsWeilDatumO}.

A continuous image of a profinite compact group in the circle is
a compact quotient of a profinite group, hence a profinite closed
subgroup of the circle.  The closed subgroups of the circle are
the finite cyclic groups and the whole circle; the latter is not
profinite.  Therefore $c$ has finite image and open kernel.
Principal finite congruences form a neighbourhood basis in this
lattice stabilizer.  Some positive integer $N$ consequently gives
\[
 \omega(k_S,k_T)v=v
       \quad\hbox{whenever }k_S,k_T\hbox{ satisfy the actual
                                      level-}N\hbox{ congruence}.
                                                               \tag{WA.9}
\]
Taking a multiple of $N$ accommodates both spaces and all lattice
denominators.  This proves the generic
\texttt{congruent\_pairs\_fix\_vacuum\_weilDatum}; it is stronger
than merely asserting that such pairs preserve a line.

\paragraph{WA5: a nonzero split-half coinvariant functional.}
Let $e$ be the actual split-place idempotent, put $\bar e=c(e)$,
and let $Y=\bar e\mathcal W(\mathbb A_f)$.  There exists a complex
linear functional $\ell$ on the ambient function space such that
\[
 \ell(v)=1,\qquad
 \ell(U(y)\phi)=\ell(\phi)quad(y\in Y, \phi\hbox{ theta}).
                                                               \tag{WA.10}
\]
No continuity of $\ell$ is needed or claimed.

Here is an algebraic proof that the required coinvariant is nonzero.
For $y\in Y\setminus\Lambda$, self-duality supplies
$\lambda\in e\Lambda$ with
$\zeta=\psi(b(\lambda,y))\ne1$.  Since $U(\lambda)v=v$,
the operator $U(\lambda)-\zeta$ kills $U(y)v$ and sends $v$ to
$(1-\zeta)v$.  Operators obtained in this way commute, since
$e\mathcal W$ is isotropic.  Multiplying one for each member of a
finite list shows that $v$ is outside the span of the translates
with $y\notin\Lambda$.  Linear extension gives a functional $q$
with
\[
                 q(U(y)v)=1_{Y\cap\Lambda}(y).              \tag{WA.11}
\]

Let $D_Y$ be the span of $U(y)\phi-\phi$ for theta functions
$\phi$ and $y\in Y$.  For each of a finite list of such differences,
smoothness at finite places gives an integer $n>0$ with
$U(ny)\phi=\phi$.  The sum
$A_y=\sum_{a=0}^{n-1}U(ay)$ kills that difference by telescoping.
All these sums commute.  Their product kills the finite list,
whereas (WA.11) evaluates its value on $v$ as a positive integer:
every summand is zero or one, and the all-zero summand is one.
Thus $v\notin D_Y$.  Take a linear functional on the quotient
which is one at $v$, and extend it to the ambient space.  This
proves (WA.10), including its invariance on every theta function.

\paragraph{WA6: make that functional algebraic on a prescribed finite list.}
For any finite list $x_1,\ldots,x_r$ with zero infinite component,
the functional in WA5 can be chosen so that
\[
                         \ell(U(x_i)v)\in\overline{\mathbb Q}
                                      \quad(1\leq i\leq r). \tag{WA.12}
\]
Write $q_0=1-e-\bar e$ and decompose
$x=ex+\bar ex+q_0x$.  If $q_0x\in\Lambda$, the invariant
$\bar e$-part and the lattice $q_0$-part can be removed using
(WA.3), with algebraic phases and a lattice sign.  For the
remaining $ex$, the value is one when $ex\in\Lambda$.  Otherwise
some $\lambda\in\bar e\Lambda$ has nontrivial commutator with
$ex$, and (WA.10) forces the value to be zero.  Thus (WA.12)
already holds in this case.

If $q_0x\notin\Lambda$, choose $u\in q_0\Lambda$ with
$\zeta=\psi(b(u,x))\ne1$, and put
$\epsilon=\varepsilon(u)$.  Replace $\ell$ by
\[
       \ell'=
       \frac{\ell\circ U(u)-\epsilon\zeta\ell}
                          {\epsilon-\epsilon\zeta}.        \tag{WA.13}
\]
The denominator is nonzero.  The replacement is still one on
$v$, and still invariant under $Y$, since $u$ is orthogonal to
$Y$.  It is zero on $U(x)v$.  For every finite $z$ its value on
$U(z)v$ is the previous value multiplied by
\[
           \frac{\epsilon\psi(b(u,z))-\epsilon\zeta}
                          {\epsilon-\epsilon\zeta}
                         \in\overline{\mathbb Q}.           \tag{WA.14}
\]
Thus already algebraic values stay algebraic.  Applying this
operation successively to the finite list proves the assertion.
The functional depends on the list, as permitted by the source;
no single algebraic functional on an infinite-dimensional space
is being assumed.

\paragraph{WA7: the splitting characters have algebraic finite values.}
For either actual splitting character $\chi$, every idele $a$
with $a_\infty=1$ satisfies $\chi(a)\in\overline{\mathbb Q}$.
The restriction condition \texttt{ResEta} makes the infinity
character trivial on positive real scalars.  Continuous circle
characters then give integers $m_\tau$ such that its infinity
part is $\prod_\tau(z_\tau/|z_\tau|)^{m_\tau}$.  At an embedded
algebraic number each factor is algebraic, because
$|z|^2=z\bar z$ is algebraic and $|z|$ is its positive square
root.  The finite integral idelic units are profinite compact;
the same circle argument as WA4 makes the character's image on
them finite, hence roots of unity.

The fractional ideal of $a$ has finite order in the ideal class
group.  Raising $a$ to that positive order and dividing by a
principal idele leaves an integral finite unit and the inverse
archimedean part of the principal idele.  Triviality on principal
ideles and the preceding two facts show $\chi(a)^h$ algebraic
for some $h>0$.  Therefore $\chi(a)$ is algebraic.  This is a
pointwise statement.  It does not put the raw values of a
half-normalized unitary character at infinitely many primes in
one number field; the source's normalized \texttt{ValuesIn}
statements retain their precise norm powers.

\paragraph{WA8: the split normalization determines the phase.}
Let $g=(g_S,g_T)$ be supported at the split pair represented by
$e+\bar e$.  Choose the exact determinant lifts $x,y$ in
\texttt{IsDetAt}.  Their infinite components are one.  The
\texttt{om\_split} clause of this actual representation says that
there is $r>0$ such that every functional as in (WA.10) satisfies
\[
 \ell(\omega(g)\phi)=r z\ell(\phi),\qquad
                   z=\chi_T(x)^{-1}\chi_S(y).               \tag{WA.15}
\]
WA7 gives $z\in\overline{\mathbb Q}$, and $\bar z z=1$.
Write $\omega(g)v=\kappa F$ as in WA3.  Unitarity and WA2 give
\[
            \bar\kappa\kappa=
                 \langle v,v\rangle/\langle F,F\rangle
                                  \in\overline{\mathbb Q}.  \tag{WA.16}
\]
Choose the functional in WA6 for the finite translates in $F$,
so $a=\ell(F)$ is algebraic.  Applying (WA.15) to $v$ gives
$\kappa a=rz$.  In particular $a\ne0$.  Taking absolute squares
and using $\bar z z=1$ yields
\[
                  r^2=(\bar\kappa\kappa)(\bar a a).
                                                               \tag{WA.17}
\]
Thus $r$ is algebraic, and $\kappa=rz/a$ is algebraic.  This
proves $\omega(g)v\in\mathcal A_v$.  The positive real scalar
and the exact character factor in (WA.15) are essential: (WA.16)
alone gives no such conclusion.

\paragraph{WA9: the algebraic-action subgroup is closed under positive roots.}
Let $H$ be the set of finite pairs $g$ for which
$\omega(g)v\in\mathcal A_v$.  If
$\omega(g)v=\kappa F$ as in WA3, covariance and WA1 give
\[
                    \omega(g)\mathcal A_v\subset\kappa\mathcal A_v.
                                                               \tag{WA.18}
\]
Indeed $\omega(g)U(x)v=U(gx)\omega(g)v$, and $gx$ is finite.
The identity and products belong to $H$.  For an inverse write
$\omega(g^{-1})v=\kappa F$.  If $g\in H$, then
$v=\kappa\omega(g)F$ with both $v$ and the nonzero
$\omega(g)F$ in $\mathcal A_v$.  WA2 makes $\kappa$ algebraic,
so $g^{-1}\in H$.

Iterating (WA.18) gives
$\omega(g^m)v=\kappa^mF_m$ with $F_m\in\mathcal A_v$ nonzero.
If $m>0$ and $g^m\in H$, WA2 makes $\kappa^m$ algebraic;
then $\kappa$ is algebraic, since it is a root of a polynomial
over a finite extension of $\mathbb Q$.  Hence
\[
                      m>0,\quad g^m\in H
                                  \quad\Longrightarrow\quad g\in H.
                                                               \tag{WA.19}
\]
WA4 puts a simultaneous principal congruence inside $H$, and
WA8 puts every split-supported finite pair inside $H$.

\paragraph{WA10: local quasi-reflections generate the remaining finite points.}
We prove directly the group statement used by the source's
\texttt{finiteUnitaryPoints\_mem\_of\_level\_and\_split\_place}.
Let $H_0$ be a subgroup of the finite unitary points which contains
a principal congruence of positive integer level $N$, contains
every split-place supported point, and satisfies positive-power
closure (WA.19).  Then $H_0$ is the whole finite unitary group.

At a nonsplit finite place, $E_w/F_v$ is a quadratic field
extension.  For a nonisotropic vector $a$, the Hermitian projection
\[
            P_a(z)=\frac{h(z,a)}{h(a,a)}a
\]
satisfies $P_a^2=P_a=P_a^*$.  If $u\bar u=1$, then
$R_{a,u}=1+(u-1)P_a$ is unitary and
\[
                       R_{a,u}^{m}=1+(u^m-1)P_a.           \tag{WA.20}
\]
These quasi-reflections generate the local unitary group.  To
verify this, take $x,z$ with the same nonzero norm.  If
$a=z-x$ is nonisotropic, put $A=h(z,a)$.  Then
$h(a,a)=A+\bar A\ne0$, so $A\ne0$, and
$u=-\bar A/A$ has norm one.  Substitution gives
$R_{a,u}z=x$.  If $z-x$ is isotropic, $z+x$ has norm
$4h(x,x)\ne0$.  The preceding construction sends $z$ to $-x$,
and $R_{x,-1}$ sends $-x$ to $x$.  A nondegenerate Hermitian
space has an orthogonal basis of nonisotropic vectors, by the
usual elimination of cross terms.  Fix its first vector using
at most two such reflections and induct on its orthogonal
complement.  This proves generation, including rank one; rank
zero is trivial.

Because conjugation preserves the local valuation,
$u\bar u=1$ gives $v_w(u)=0$.  Every finite quotient
$\mathcal O_{E_w}^{\times}/(1+\mathfrak p_w^b)$ is finite.
Some positive power $u^m$ is therefore congruent to one to any
prescribed depth.  Choose that depth large enough to clear the
matrix denominators of $P_a$ relative to the fixed local lattice
and to meet level $N$.  Formula (WA.20) then puts $R_{a,u}^m$
in the prescribed principal congruence; its inverse is integral
and has the same congruence by the convergent geometric-series
inverse of a sufficiently small integral matrix perturbation.
Power closure puts $R_{a,u}$ itself in $H_0$.  No condition that
$m$ be prime to the residue characteristic is needed.  This
argument therefore includes ramified and dyadic places.

Finally any finite adelic point stabilizes the fixed lattice
away from a finite set of places.  Enlarge that set by the places
dividing $N$.  Removing its finitely many single-place factors
leaves a member of the principal congruence, since $N$ is a
local unit everywhere else.  Split factors are already in
$H_0$; nonsplit factors are products of the reflections just
treated.  Their product proves the assertion.  This local proof
replaces the source's longer rational/split/congruence
factorization, without changing its group-theoretic conclusion.

\paragraph{WA11: actual finite Weil operators preserve the algebraic span.}
Apply WA10 to the inverse image of the subgroup in WA9 under
$g_S\mapsto(g_S,1)$ and then under $g_T\mapsto(1,g_T)$.
All its hypotheses were proved in WA4, WA8 and WA9.  Thus, for
every actual finite pair,
\[
     \omega(g_S,g_T)v\in\mathcal A_v,
          \qquad \omega(g_S,g_T)\mathcal A_v=\mathcal A_v.   \tag{WA.21}
\]
The equality follows from covariance and the same assertion for
the inverse.  This is the literal conclusion of
\texttt{BotExpand.om\_vac\_algComb}.  The source's
\texttt{moved\_translate\_algComb} follows by applying it to
$U(x)v$ and using (WA.3), with every resulting translation still
having infinite component zero.

The proof uses the displayed generic data; it does not use the
nominal \texttt{FunPr} or universal \texttt{NormSum} arguments
carried by the bottom wrappers.  It proves qualitative
algebraicity of finite coefficients.  It does not prove that an
arbitrary distinguished dictionary scalar has absolute value
one under every algebraic embedding, or is a unit at specified
finite primes.  Those stronger assertions in
\texttt{ExistsUnitModulusScalarLocal} remain distinct.

\paragraph{WA12: both actual bottom datum functions are algebraic combinations.}
Fix the finite half-space and the completed half-lattice in a
literal \texttt{IsBottomTri}.  Each of its two standard finite
functions is algebraic-valued.  Its determinant character is
evaluated on an integral finite unit, hence has finite order by
the profinite argument of WA7; its remaining factors are
indicators.  On a suitable bounded fractional half-lattice the
function is periodic by $k$ times the integral half-lattice,
for one positive integer $k$.  It is consequently a finite
linear combination of small-coset indicators.

For one such indicator, finite character orthogonality gives
\[
  1_{a+L_0}(x)=\frac1{|L/L_0|}
        \sum_{\alpha\in(L/L_0)^\vee}\alpha(x-a)
                                      \quad (x\in L),       \tag{WA.22}
\]
where $L$ contains the bounded support and $L_0$ is a common
period lattice.  The dual half-lattice realizes these characters
by modulation.  Translation and modulation are the two
isotropic Heisenberg actions; converting their product to one
$U(w)$ introduces exactly the algebraic phase of WA1.
The actual theta map is linear, sends the prescribed basic
half-lattice indicator to the given vacuum, and intertwines
these translations and modulations.  Thus (WA.22) proves
\[
             X.f_t\in\mathcal A_{X.\mathrm{vac}},\qquad
             X.f'_t\in\mathcal A_{X.\mathrm{vac}'}.          \tag{WA.23}
\]
On the second side use the source's conjugate idempotent and its
own actual character and vacuum.  No Fourier transform is
silently substituted for that second standard function.

WA11 applies separately to the two actual oriented Weil data,
so (WA.23) remains true after every finite pair acts.  Finite
Heisenberg operators commute with archimedean polynomial
creation, by covariance at disjoint places and the $L^2$
difference-quotient justification in WA3.  Therefore every
monomial bottom kernel is a finite algebraic sum of finite
translate jets of its own vacuum.  This proves the actual
\texttt{theta\_kernel\_algebraic\_expansion}, its dual, and
\texttt{ker'\_eq\_sum\_jets} with their original orientations.

\paragraph{WA13: the actual normalized grid jets are algebraic.}
Fix one literal \texttt{Choice}, its grid Weil representations,
normal datum and rational frames.  Every edge used in TT13 has
the self-dual lattices and actual unit vacuum of TT12, with the
specified Cartan structures.  WA11 supplies (TT.19) for this
very representation.  For each pure coset datum, the same
finite Fourier calculation as WA12 supplies its vector as an
algebraic finite sum of translates of that vacuum.

The period model of that normal datum is the one already used
in TA10 and VI9--VI11.  Finite additive approximation identifies
a finite-adelic translate with a rational archimedean translate,
up to the lattice law and a root of unity from WA1.  Thus its
monomial jet has the exact form
\[
 c_{0,p}\prod_\tau
     \operatorname{opp}_p(\operatorname{period}(\delta,\tau))
                         ^{d_\tau}\,a,
                      \qquad a\in\overline{\mathbb Q}.      \tag{WA.24}
\]
Here $c_{0,p}$ is the actual nonzero comparison scalar of the normal datum;
it is not asserted algebraic.  The half-point, $\delta$, vacuum
and frame in (WA.24) are unchanged.  TT13 proves that the
homogeneous degree of the actual differential operator factors
out exactly the period powers occurring in
$Z.\mathrm{scal}(1,p)$, whose definition is their product with
$c_{0,p}$ inverted.  They are nonzero and cancel.  TT9's rational
polynomial and bideterminant operations then give
\[
                  Z.\mathrm{kk}(d,x,y)_{ab}
                                \in\overline{\mathbb Q}     \tag{WA.25}
\]
for every actual required entry.  This supplies the previously
retained arithmetic premise of TT13.  It uses the direct
rational-translate germ theorem, not a claim that arbitrary
algebraic entire functions are stable under translation.

\paragraph{WA14: the literal normalized jet field now exists.}
For a fixed choice, TT6--TT8 give finite support labels and
finite class sets.  At the least depth for each label, TT7
shows that the support of every allowed expansion is contained
in one finite set.  Taking all twists, top and row edges and
matrix indices therefore gives a genuinely finite set
$E_{\rm jet}$ of all entries demanded by \texttt{JetsIn}.
WA13 makes every member algebraic.  Consequently
\[
                   M_{\rm jet}=\mathbb Q(E_{\rm jet})
\]
is a number field, and its inclusion into $\mathbb C$ satisfies
the literal \texttt{JetCl} for the original choice.  Combining
this with TT8 gives \texttt{NFJet} with no new size threshold:
the two additional thresholds may be taken to be zero.  This
is a qualitative field-existence result; it gives no uniform
degree, height or denominator bound.  In particular it does
not replace the quantitative clauses of \texttt{DictAt}.

\paragraph{WA15: one fixed-degree period for the actual bottom vacua.}
Let $X$ be a literal bottom triangle and let $d$ be a fixed
place-degree vector.  Its positive spaces have standard Gram
matrices, integral lattices and identity orthonormal frames.
The exact square-rank identities make the total target rank
$\operatorname{rank}WW=\operatorname{botBig}+\operatorname{botSmall}$
even.  Hence both products
$\operatorname{rank}A\operatorname{rank}WW$ and
$\operatorname{rank}A'\operatorname{rank}WW$ are even.  We use
this product parity, not an assertion that both source ranks
are even.

Conjugating the negative source vacuum gives a positive vacuum
for $A,WW$.  Choose the same standard CM curve datum and half-point
as in the earlier theta construction.  The positive period-model
function and this actual vacuum satisfy the same lattice and Fock
laws and are both nonzero.  Analytic uniqueness from WA3 therefore
gives one constant $c\in\mathbb C^{\times}$ relating them.  The
period-normalized entire function need not have unit $L^2$ norm,
so no absolute-value normalization of $c$ is asserted.
VI9--VI11 and the exact Fock derivative formula then give the
common nonzero period
\[
          \Omega_d=\overline{c\prod_\tau
                      \operatorname{period}(\delta,\tau)^{d_\tau}}
                                                               \tag{WA.26}
\]
for all finite translates and all monomials of that place degree.
The negative-space identity conjugates the whole positive jet,
including $c$ and the period product.  On the positive source
$A'$ use its own nonzero comparison scalar $c'$ and the unconjugated
product.  All periods are nonzero; degree zero and rank zero
cause no division by a zero quantity.

WA12 expresses every finite-point kernel as an algebraic sum
of precisely these translate jets.  Division by (WA.26) therefore
makes every monomial kernel of degree $d$ algebraic.  This proves
the actual \texttt{period\_algebraic\_fixed\_degree} and both
finite-point kernel arithmetic statements without an independent
\texttt{NormSum.alg} or finite-Weil arithmetic premise.  The
comparison scalars are absorbed in the periods; no algebraicity
of $c,c'$ is needed.

\paragraph{WA16: the actual definite integrals have a common number field.}
For every supplied bottom triangle $X$, there is a number field
$M\subset\mathbb C$ and nonzero periods $\Pi,\Pi'$ such that
the literal \texttt{SumsIn} holds.  We now discharge TT16's
remaining finite-Weil premise by WA11--WA12 and its fixed-degree
period input by WA15.

The argument retains the actual measures.  A common compact open
finite level fixes the finitely many standard theta data, using
WA4 and the finite Heisenberg expansions.  Polynomial creation
commutes with that level, so it is independent of the monomial.
The definite adelic quotient is compact and has finitely many
classes at this level.  Its exact integral decomposition is
\[
 I_t(P,h)=v_0\sum_l|\Gamma_l|^{-1}
       \int_C\operatorname{ker}_t(P,h)(z_lk)\,dk,
                 \qquad v_0>0,                             \tag{WA.27}
\]
with probability Haar measure on the compact real group $C$.
The constant $v_0$ is the original compact-open volume; it need
not be rational.  The continuous kernels are integrable on this
compact quotient.  TT15 identifies the compact projector over
$\overline{\mathbb Q}$ and the single surviving degree $d_0$,
using the exact determinant weight.  WA15 then proves
$I_t(P,h)/(v_0\Omega_{d_0})\in\overline{\mathbb Q}$ for all
finite $h$ and monomials $P$.  The statement also holds when all
these integrals are zero; the chosen period remains nonzero.

To obtain one finite field, rather than infinitely many algebraic
values, use the finite generation proved in TT16.  A common
finite level and rational invariance reduce the finite points to
finitely many class representatives.  The compensating infinite
matrices are embeddings of rational $E$-matrices.  Fock
substitution and integral determinant powers have coefficients
in one fixed finite normal extension of $E$.  The central
character equation permits only finitely many degree vectors
and finitely many monomials in them.  Thus every integral is a
linear combination over one number field $M_0$ of a finite list
$I_t(P_b,h_a)$.  Adjoin their normalized values, already known
algebraic.  Repeat the argument for $I'_t$ with its own measure
and period.  The compositum is finite and supplies exactly
\texttt{SumsIn}, with the same $X$ on both sides of the join.

\paragraph{WA17: the standard bottom spaces and the precise remaining existence.}
For each square $(i,j)$ the source defines the two nonnegative
ranks \texttt{botBig}, \texttt{botSmall}.  Take $A=E^{\rm botBig}$,
$A'=E^{\rm botSmall}$ with Gram matrix one and lattice
$\mathcal O_E^n$, and take $WW=A\oplus A'$ with those same
standard choices.  Testing the Hermitian dual against the
coordinate vectors shows that each standard lattice is
self-dual.  The Cartan involutions and orthonormal frames are
identities; the direct-sum rank identity and the product parity
in WA15 follow from the literal square-rank formulas.

The bottom characters are the source's parity corrections of
$\mu_t(j+1)$ and $\chi_{m-i}$.  Multiplication by the prescribed
power $\eta^{\operatorname{rank}V\bmod2}$ gives their exact
\texttt{ResEta} ranks; this is addition of the integral
restriction exponents modulo two.  It does not permit changing
the finite-order twist or choosing a new $\eta$ independently.
Given the two full oriented Weil representations for
$(\operatorname{neg}A,WW)$ and $(A',WW)$ with these characters,
choose the normalized vacua, the original finite split
idempotent, the actual conductor exponents, and the two theta
maps for the standard finite functions.  Choose Haar measures
and fundamental domains for these actual groups.  These choices
give the literal bottom triangle: its two rigs, target rig,
representations, vacua, idempotent, conductors, finite functions,
measures and domains have exactly the stipulated identities.
WA16 then gives its period field.  A finite compositum over all
squares supplies \texttt{NFBot} under this exact existence input.

The new representations here are not the already supplied grid
representations.  Their full existence is not proved in
WA1--WA16.  The next paragraph narrows it to a canonical
splitting-character case by an explicit construction.

\paragraph{WA18: change splitting characters by a determinant-root character.}
Suppose a full oriented Weil datum exists for the same $S,T$
with characters $\eta^{\operatorname{rank}S}$ and
$\eta^{\operatorname{rank}T}$, where $\eta$ is the same
unramified splitting character.  Then a full datum exists for
any of the prescribed characters $\chi_S,\chi_T$ with the
required rank restrictions.  Put
\[
 \nu_S=\eta^{-\operatorname{rank}S}\chi_S,\qquad
 \nu_T=\eta^{-\operatorname{rank}T}\chi_T.
                                                               \tag{WA.28}
\]
They are Hecke characters trivial on the ideles of $F$.
For a unitary element $g$, choose a determinant root $x$ with
$x/\bar x=\det g$.  Hilbert 90 supplies such roots locally.
Outside a finite set the extension and $\nu$ are unramified.
There a norm-one unit has a unit root: in the quadratic field
case rescale any root by the appropriate power of a base-field
uniformizer, whose extension valuation is one; in the split
case use the explicit root $(t,1)$.  These roots assemble to
an idele.  Moreover $\widetilde\nu$ is trivial on the integral
unitary group at every such place, because $\nu$ is trivial on
the local integral units.
Two choices differ by an $F$-idele, so $\nu(x)$ is independent
of the choice.  This defines a continuous unitary character
$\widetilde\nu(g)$.  Multiplicativity follows by using the
product of two roots; rational unitary elements have rational
roots by the field form of Hilbert 90, so this character is
trivial on rational points.  Continuity on the restricted product
now reduces to the finitely many remaining factors.  There it
follows from local
sections of $x\mapsto x/\bar x$: near one solve the equation
by the nonzero expression $1+t$ for $t\bar t=1$, with a
fixed alternative chart at $t=-1$.

On the same theta functions set
\[
 \omega'(g,h)=\widetilde\nu_T(g)^{-1}
                          \widetilde\nu_S(h)\,\omega(g,h).
                                                               \tag{WA.29}
\]
The two characters are deliberately crossed.  Their scalars
commute with Heisenberg operators, have absolute value one,
are continuous, and are trivial at rational pairs.  Therefore
representation, covariance, automorphy, continuity and norm
preservation all hold.  At a split pair the determinant roots
can be the literal \texttt{IsDetAt} lifts, so the scalar in
(WA.29) changes the split factor from
$\eta(x)^{-\operatorname{rank}T}\eta(y)^{\operatorname{rank}S}$
to $\chi_T(x)^{-1}\chi_S(y)$, with the same positive real
multiplier.  This verifies the last defining clause, rather than
merely asserting that a twist is harmless.

Thus the remaining full-Weil supplier can be taken to be
\texttt{exists\_weilDatum\_etaTilde}: the actual construction
of its finite and real actions, their commutation and rational
theta-distribution automorphy.  The source constructs these;
their reconstruction is a separate target.  Nothing in this
section treats that endpoint as an axiom-free consequence of
the character adjustment.

\paragraph{WA19: the same-choice field consequence and its exact boundary.}
For a fixed literal choice, WA14 now supplies its jet field.
For each supplied bottom triangle, WA16 supplies its two-period
field.  Taking their finite compositum with the character,
root and phase fields of TT11 gives these portions of
\texttt{NumbersIn} on the same choice and the same embedded
complex realization.  If bottom Weil data are supplied through
WA17--WA18 and the canonical case retained there, the finite
compositum covers all squares as demanded by \texttt{NFBot}.

Critical toric $L$-value ratios are not finite Heisenberg
coefficients.  Their remaining source input is
\texttt{ToricRatioConjugates} with its actual CM field data
and reference twist.  A field containing those ratios can be
joined with the fields above, but its existence is not inferred
here from (WA.21).  Likewise, membership in a subfield of
$\mathbb C$ is distinct from membership in the range of the
already chosen embedding $\sigma:K\hookrightarrow\mathbb C$.
The literal \texttt{Choice.FieldOK} still requires its same
normalization field, frame, Hida setting, family and specified
embedding to agree.  TT19's compatibility argument applies
once those actual joins and the remaining ratio field are
supplied.  No replacement choice or independent collection of
embeddings is silently substituted for them.

The substantive new endpoints are consequently (WA.21),
\texttt{NFJet}, and \texttt{SumsIn} for the actual supplied
bottom data.  The quantitative all-embedding and finite-unit
dictionary estimates, the canonical full bottom-Weil existence,
the critical ratio field, and the remaining compatible-choice
join retain their stated source scopes.
